\documentclass[11pt]{article}
\usepackage[margin=0.78in]{geometry}
\usepackage{amsmath,amssymb,amsthm,mathtools}
\usepackage{graphicx}
\usepackage{booktabs}
\usepackage{array}
\usepackage{enumitem}
\usepackage{hyperref}
\usepackage{xcolor}
\usepackage{microtype}
\usepackage{fancyhdr}
\usepackage{listings}
\usepackage{caption}
\hypersetup{colorlinks=true,linkcolor=blue!45!black,citecolor=blue!45!black,urlcolor=blue!45!black}
\pagestyle{fancy}
\fancyhf{}
\lhead{Geometric Infrastructure Optimization v2}
\rhead{Hard-Push Research Bundle}
\cfoot{\thepage}

\newtheorem{definition}{Definition}[section]
\newtheorem{theorem}{Theorem}[section]
\newtheorem{proposition}{Proposition}[section]
\newtheorem{lemma}{Lemma}[section]
\newtheorem{corollary}{Corollary}[section]
\theoremstyle{remark}
\newtheorem{remark}{Remark}[section]

\lstset{basicstyle=\ttfamily\small,breaklines=true,frame=single,columns=fullflexible,showstringspaces=false}

\title{\vspace{-1.0cm}Geometric Infrastructure Optimization v2\\
\large Locality-First Compute Fabrics, Exact Swap Auditing, and Machine-Precision Tests}
\author{Degen Studios / BoredNBased Research Draft}
\date{May 28, 2026}

\begin{document}
\maketitle

\begin{abstract}
Parallelism increases concurrency; geometry determines the communication tax paid by that concurrency. This v2 hard-push formalizes a locality-first infrastructure layer that embeds a workload communication graph into a physical server, rack, switch, row, pod, or optical fabric metric. The goal is not to replace Clos, fat-tree, dragonfly, or leaf-spine networks; the goal is to add a live placement layer that makes hot logical edges physically short whenever the workload graph contains locality. The companion notebook verifies the exact swap-delta identity, monotone descent, exhaustive small-case optimum, float64 exactness inside the $2^{53}$ envelope, space-filling/ribbon locality, rack-aware cluster placement, dynamic drift with migration cost, and null cases where geometry should not help. The result is a falsifiable scheduler theory: it wins when communication has locality, hierarchy, or hot communities, and it correctly collapses to no advantage under uniform communication.
\end{abstract}

\paragraph{Executed evidence root.}
\[
\texttt{cff9b39527b2746771d855214e3d7e1d7aa3092bfb990aa7854c9393363afd2c}
\]
This root is the SHA-256 digest of the canonical evidence JSON produced by the executed notebook. It authenticates the test manifest, not physical-hardware performance.

\section{One-line thesis}

A compute fabric should be treated as a geometric substrate. The scheduler should continuously embed the workload communication graph into the physical network geometry so high-weight logical edges occupy low-cost physical distances. In operational terms:
\[
\text{parallel lanes} + \text{geometric placement} + \text{telemetry feedback} = \text{lower communication tax.}
\]

\section{Position relative to known infrastructure}

The proposal is compatible with existing topology work. Fat-tree and Clos-like data-center networks provide high aggregate capacity and multiple paths between hosts \cite{alfares2008fat,singh2015jupiter}. Dragonfly networks reduce global diameter through high-radix local groups and sparse global links \cite{kim2008dragonfly}. VL2 and related designs decouple addressing from location and use path diversity for scale \cite{greenberg2009vl2}. Traffic-aware VM placement work already shows that workload placement and network traffic are coupled, and that placement can reduce network cost \cite{meng2010improving,jiang2012joint}. Space-filling curves, especially Hilbert-style orderings, are well known locality-preserving mappings for transforming multi-dimensional neighborhoods into one-dimensional traversals \cite{moon2001hilbert,mokbel2003space}.

This document adds a focused infrastructure layer:
\begin{quote}
Topology gives the possible routes. Geometric placement changes which routes matter.
\end{quote}

\section{Formal model}

\begin{definition}[Logical communication graph]
Let $G=(V,E,W)$ be an undirected weighted graph over $n$ logical workloads. The nonnegative integer matrix $W\in\mathbb{Z}_{\ge 0}^{n\times n}$ is symmetric with $W_{ii}=0$. A larger $W_{ij}$ means workloads $i$ and $j$ communicate more frequently, move more bytes, synchronize more often, or have higher latency sensitivity.
\end{definition}

\begin{definition}[Physical fabric metric]
Let $P=\{0,\ldots,n-1\}$ be physical slots, such as GPU slots, hosts, rack units, server nodes, leaf-switch neighborhoods, rows, pods, or sites. The symmetric matrix $D\in\mathbb{Z}_{\ge 0}^{n\times n}$ satisfies $D_{aa}=0$ and encodes physical communication cost. This cost can represent hop count, cable distance, optics, congestion, serialization, power, cooling zone penalty, failure-domain penalty, or a weighted blend of these.
\end{definition}

\begin{definition}[Placement]
A placement is a permutation $\pi\in S_n$ where $\pi_i$ is the physical slot assigned to logical workload $i$.
\end{definition}

\begin{definition}[Weighted geometric communication cost]
\begin{equation}
J(\pi)=\sum_{0\le i<j<n} W_{ij}D_{\pi_i,\pi_j}.
\label{eq:objective}
\end{equation}
The scheduler's idealized objective is to minimize $J$ subject to operational constraints.
\end{definition}

\paragraph{Interpretation.} Equation \eqref{eq:objective} says that a heavily communicating pair is expensive only when the two workloads are physically far apart. A lightly communicating pair can be far apart with lower penalty. The scheduler is therefore solving a weighted graph embedding / quadratic assignment style problem.

\section{Exact swap-delta theorem}

A live scheduler cannot recompute an entire $n^2$ score for every possible small move. It needs a cheap exact test for local moves.

\begin{theorem}[Exact two-swap delta]
Let $\pi'$ be obtained from $\pi$ by swapping the physical assignments of logical workloads $a$ and $b$. Then
\begin{equation}
J(\pi')-J(\pi)=\sum_{k\ne a,b}(W_{ak}-W_{bk})\left(D_{\pi_b,\pi_k}-D_{\pi_a,\pi_k}\right).
\label{eq:swapdelta}
\end{equation}
\end{theorem}

\begin{proof}
All pair contributions not incident to $a$ or $b$ are unchanged. The edge $(a,b)$ is also unchanged because $D$ is symmetric:
$D_{\pi_b,\pi_a}=D_{\pi_a,\pi_b}$. For each $k\ne a,b$, the old incident contribution is
\[
W_{ak}D_{\pi_a,\pi_k}+W_{bk}D_{\pi_b,\pi_k}.
\]
The new contribution is
\[
W_{ak}D_{\pi_b,\pi_k}+W_{bk}D_{\pi_a,\pi_k}.
\]
Subtracting old from new yields
\[
W_{ak}(D_{\pi_b,\pi_k}-D_{\pi_a,\pi_k})+W_{bk}(D_{\pi_a,\pi_k}-D_{\pi_b,\pi_k}),
\]
which simplifies to Equation \eqref{eq:swapdelta}. Summing over all $k\ne a,b$ completes the proof.
\end{proof}

\paragraph{Operational meaning.} A rack-aware scheduler can evaluate a candidate swap in $O(n)$ integer operations instead of rescoring the entire fabric. Every accepted swap can be audited by verifying
\[
J(\pi')-J(\pi)=\Delta_{a,b}.
\]
The v2 notebook tests this identity over 300 randomized permutations and 3000 sampled swaps with zero failures.

\section{Local descent and termination}

\begin{definition}[Migration-aware local move]
Let $M_{a,b}\ge 0$ be the operational migration cost of swapping workloads $a$ and $b$. The swap is accepted only if
\[
\Delta_{a,b}+M_{a,b}<0.
\]
\end{definition}

\begin{proposition}[Finite termination]
If $W$ and $D$ are integer matrices and the algorithm accepts only swaps with $\Delta_{a,b}+M_{a,b}<0$, then the unadjusted cost $J$ strictly decreases for every accepted zero-migration move and the adjusted objective decreases for every positive-migration move. Since there are finitely many placements, the process terminates at a one-swap local optimum.
\end{proposition}

\begin{remark}
This is not a guarantee of global optimality. The global problem is a quadratic assignment style combinatorial optimization problem. The practical win is that the local optimizer is exact, monotone, auditable, and can be used as a safe real-time improvement layer.
\end{remark}

\section{Machine-precision discipline}

The v2 harness keeps the core placement objective in integer arithmetic. It then separately checks float64 display equivalence inside the exact-integer range of IEEE-754 double precision.

\begin{proposition}[Safe integer display envelope]
If $J(\pi)<2^{53}$ and $J(\pi)$ is an integer, then converting $J(\pi)$ to float64 and back preserves the integer exactly. The executed notebook tests 500 deterministic placements with maximum integer objective $49168$ and maximum float error $0.0$.
\end{proposition}

\section{When geometry wins}

The objective has leverage when $W$ is not uniform. The strongest cases are:
\begin{enumerate}[leftmargin=*]
\item \textbf{Clustered AI services:} routers, model shards, KV caches, vector shards, and rerankers form hot neighborhoods.
\item \textbf{Distributed training:} gradient synchronization, pipeline stages, tensor-parallel groups, and parameter shards have repeated communication patterns.
\item \textbf{Simulation cells:} boundary conditions imply nearest-neighbor traffic in logical space.
\item \textbf{Storage systems:} hot objects, metadata, indexes, and related vector/document shards benefit from physical co-location.
\item \textbf{Agent colonies:} local agents should exchange dense state locally and export only summarized state globally.
\end{enumerate}

The v2 rack experiment uses a four-community communication graph over 32 logical workloads and a rack metric over four physical rack neighborhoods. A random placement scored $3{,}492{,}265$. The local optimizer reached $810{,}605$, and deterministic cluster packing scored $811{,}270$. That is a $76.77\%$ reduction from the random start in this synthetic instance.

\section{When geometry should fail}

A good theory must define its null case.

\begin{theorem}[Uniform complete communication null case]
If $W_{ij}=c$ for all $i\ne j$ and $D$ is fixed, then $J(\pi)$ is the same for every placement $\pi$.
\end{theorem}

\begin{proof}
For any permutation $\pi$,
\[
J(\pi)=c\sum_{i<j}D_{\pi_i,\pi_j}.
\]
Because $\pi$ merely reorders the physical slots, the unordered set of pairs $\{\{\pi_i,\pi_j\}:i<j\}$ is exactly the unordered set of all physical slot pairs. The sum is therefore invariant.
\end{proof}

\paragraph{Falsification result.} The v2 notebook confirms this exactly: 50 random permutations of a uniform complete communication matrix all produce the same objective. This matters because the theory does not pretend geometry helps when there is no locality to exploit.

\section{Space-filling/ribbon placement}

When a two-dimensional logical manifold must be placed along a one-dimensional physical ribbon, aisle, rack sequence, or cable run, space-filling curves become useful heuristics. The notebook compares Hilbert, Morton/Z-order, snake, row-major, and random placements over an $8\times 8$ logical lattice with distance-decaying communication.

\begin{table}[h]
\centering
\caption{Ribbon-placement cost comparison from executed v2 notebook. Lower is better.}
\begin{tabular}{lr}
\toprule
Placement & Objective $J$ \\
\midrule
Hilbert & $1{,}050{,}248$ \\
Morton & $992{,}408$ \\
Snake & $971{,}568$ \\
Row-major & $971{,}568$ \\
Random mean & $2{,}317{,}939.02$ \\
\bottomrule
\end{tabular}
\end{table}

Hilbert was $54.69\%$ lower than random mean in this run. The best named deterministic ribbon order was snake/row-major for this particular synthetic lattice. That is an important honest result: Hilbert is a strong locality-preserving family, but it is not universally dominant. The scheduler should benchmark candidate embeddings against the actual $W$ and $D$ rather than assume one curve is always best.

\begin{figure}[h]
\centering
\includegraphics[width=0.82\linewidth]{geo_v2_figures/ribbon_order_comparison.png}
\caption{Logical locality mapped onto a one-dimensional physical ribbon.}
\end{figure}

\section{Dynamic drift and anti-thrash migration cost}

Real workloads drift. A good system must avoid migration thrash. The v2 dynamic test changes $W$ after initial placement and accepts a remap only when the exact integer delta clears a migration threshold. The test result:
\begin{itemize}[leftmargin=*]
\item Cost on drifted workload before remap: $2{,}872{,}404$.
\item Cost after thresholded remap: $592{,}521$.
\item Accepted remap moves: $9$.
\item Migration cost per swap: $1000$.
\item Drift reduction: $79.37\%$.
\end{itemize}

\begin{figure}[h]
\centering
\includegraphics[width=0.82\linewidth]{geo_v2_figures/dynamic_drift_remap.png}
\caption{Thresholded remap under workload drift. Every accepted move clears migration cost.}
\end{figure}

\section{Live scheduler architecture}

The production architecture has five loops.

\subsection{Telemetry loop}
\begin{lstlisting}
collect: bytes/sec, packets/sec, RPC count, tail latency,
         queue depth, all-reduce frequency, cache hits,
         vector shard calls, model-router calls, storage affinity
emit:    edge observations (i, j, weight, confidence, timestamp)
\end{lstlisting}

\subsection{Graph-building loop}
\begin{lstlisting}
W_t = exponential_decay(W_{t-1}) + new_edge_observations
W_t = denoise(W_t, confidence, service class)
W_t = cap_or_transform(W_t)       # prevent one metric from dominating
\end{lstlisting}

\subsection{Physical metric loop}
\begin{lstlisting}
D = weighted_sum(
      hop_count,
      switch_stage_penalty,
      optical_conversion_penalty,
      cable_distance,
      congestion_shadow_price,
      thermal_zone_penalty,
      failure_domain_penalty,
      policy_boundary_penalty
    )
\end{lstlisting}

\subsection{Optimizer loop}
\begin{lstlisting}
for candidate swap (a,b):
    delta = exact_swap_delta(W, D, placement, a, b)
    adjusted = delta + migration_cost(a,b)
    if adjusted < 0 and safety_constraints_pass:
        stage move
simulate staged moves
commit only if total predicted gain clears threshold
\end{lstlisting}

\subsection{Verification loop}
\begin{lstlisting}
measure before/after latency, traffic, hops, power proxies
compare observed gain to predicted gain
update confidence model
rollback or quarantine if regression exceeds guardrail
\end{lstlisting}

\section{Result manifest}

\begin{table}[h]
\centering
\caption{Executed v2 tests.}
\begin{tabular}{lll}
\toprule
Test & Result & Key value \\
\midrule
Swap-delta identity & PASS & 300 trials, zero failures \\
Float64 exactness & PASS & max error $0.0$ \\
Exhaustive small optimum & PASS & best gap $0$ \\
Ribbon locality & PASS & Hilbert $54.69\%$ below random mean \\
Rack cluster placement & PASS & $76.77\%$ reduction from random \\
Dynamic drift & PASS & $79.37\%$ reduction after remap \\
Uniform null case & PASS & invariant objective \\
Expander-like near-null & PASS & only $3.47\%$ local reduction \\
\bottomrule
\end{tabular}
\end{table}

\begin{figure}[h]
\centering
\includegraphics[width=0.82\linewidth]{geo_v2_figures/swap_descent_trajectory.png}
\caption{Monotone local swap descent from a random placement.}
\end{figure}

\begin{figure}[h]
\centering
\includegraphics[width=0.82\linewidth]{geo_v2_figures/rack_cluster_comparison.png}
\caption{Rack-aware cluster placement versus random start and local descent.}
\end{figure}

\section{Practical deployment target}

A real prototype can start without changing rack hardware. The first version only observes and recommends:
\begin{enumerate}[leftmargin=*]
\item Build $W$ from traffic telemetry, service calls, all-reduce traces, storage access, or application-level spans.
\item Build $D$ from inventory: host, rack, row, leaf, spine, pod, power, cooling, and policy metadata.
\item Score current placement.
\item Generate safe candidate migrations.
\item Simulate cost savings and operational risk.
\item Apply only to low-risk services first.
\item Compare predicted savings to observed latency/power/hop improvements.
\end{enumerate}

The same mechanism can be applied at multiple scales: GPU slots inside a host, hosts inside a rack, racks inside a row, pods inside a region, or federated browser/edge nodes across a shop floor.

\section{Conclusion}

The v2 result is stronger than the initial statement. The correct claim is not that geometric layout always doubles optimization. The precise claim is:
\begin{quote}
If a workload communication graph has locality, hierarchy, hot communities, or repeated nearest-neighbor structure, then embedding that graph into physical network geometry can substantially reduce weighted communication distance. The improvement is exactly scoreable, locally auditable, and falsifiable under uniform communication.
\end{quote}

The machine-precision tests support the mathematical core. Hardware validation is the next step: collect real $W$, build a real $D$, run the scheduler in observe-only mode, and compare predicted versus measured reductions.

\begin{thebibliography}{9}

\bibitem{alfares2008fat}
M. Al-Fares, A. Loukissas, and A. Vahdat.
\newblock A scalable, commodity data center network architecture.
\newblock \emph{ACM SIGCOMM}, 2008.
\newblock \url{https://ccr.sigcomm.org/online/files/p63-alfares.pdf}

\bibitem{kim2008dragonfly}
J. Kim, W. J. Dally, S. Scott, and D. Abts.
\newblock Technology-driven, highly-scalable dragonfly topology.
\newblock \emph{ISCA}, 2008.
\newblock \url{https://research.google.com/pubs/archive/34926.pdf}

\bibitem{greenberg2009vl2}
A. Greenberg et al.
\newblock VL2: A scalable and flexible data center network.
\newblock \emph{ACM SIGCOMM}, 2009.

\bibitem{singh2015jupiter}
A. Singh et al.
\newblock Jupiter rising: A decade of Clos topologies and centralized control in Google's datacenter network.
\newblock \emph{ACM SIGCOMM}, 2015.

\bibitem{meng2010improving}
X. Meng, V. Pappas, and L. Zhang.
\newblock Improving the scalability of data center networks with traffic-aware virtual machine placement.
\newblock \emph{IEEE INFOCOM}, 2010.

\bibitem{jiang2012joint}
J. W. Jiang et al.
\newblock Joint VM placement and routing for data center traffic engineering.
\newblock \emph{IEEE INFOCOM}, 2012.

\bibitem{moon2001hilbert}
B. Moon, H. V. Jagadish, C. Faloutsos, and J. H. Saltz.
\newblock Analysis of the clustering properties of the Hilbert space-filling curve.
\newblock \emph{IEEE Transactions on Knowledge and Data Engineering}, 2001.
\newblock \url{https://www.cs.cmu.edu/~christos/PUBLICATIONS/ieee-tkde-hilbert.pdf}

\bibitem{mokbel2003space}
M. F. Mokbel, W. G. Aref, and I. Kamel.
\newblock Analysis of multi-dimensional space-filling curves.
\newblock \emph{GeoInformatica}, 2003.

\end{thebibliography}

\end{document}
